\chapter{Estructura Organizativa de una Planta Industrial de Producción de Stevia}

\section{Introducción y Contexto del Proyecto}

El presente caso práctico aborda el diseño y definición de la \textbf{estructura organizativa y de recursos humanos} para una nueva planta industrial orientada a la extracción, purificación y cristalización de \textbf{glucósidos de esteviol} (principalmente esteviósido y rebaudiósido A) a partir de hoja seca de \textit{Stevia rebaudiana Bertoni}.

Al tratarse del diseño de una planta dentro de la categoría de \textbf{pequeña/mediana empresa industrial (PYME)}, se adopta una plantilla dimensionada en torno a \textbf{31 trabajadores}, optimizando la polivalencia, la cobertura de turnos en producción continua y el estricto cumplimiento de normativas de calidad y seguridad alimentaria (normas IFS/BRC, ISO 22000 y GMP).

\begin{note}
    \textbf{Aspectos clave de la planta de Stevia:}
    \begin{itemize}
        \item \textbf{Proceso productivo:} Recepción y molienda de hoja seca $\rightarrow$ Extracción sólido-líquido (acuosa/etanol) $\rightarrow$ Coagulación/Floculación $\rightarrow$ Filtración por membranas/resinas de intercambio iónico y adsorción $\rightarrow$ Cristalización fraccionada $\rightarrow$ Secado por atomización/lecho fluido $\rightarrow$ Envasado aséptico.
        \item \textbf{Sector:} Ingrediente alimentario de alto valor añadido y pureza ($\geq 95\%$ glucósidos de esteviol).
        \item \textbf{Estructura organizativa adoptada:} Funcional-jerárquica con coordinación matricial por proyectos y turnos rotativos en planta.
    \end{itemize}
\end{note}

\section{Dimensionamiento Global de la Plantilla}

La empresa contará con un total de \textbf{31 personas}, distribuidas en 7 departamentos funcionales:

\begin{table}[H]
    \centering
    \caption{Distribución global de personal por departamentos}
    \label{tab:resumen_personal}
    \begin{tabular}{@{}llcc@{}}
        \toprule
        \textbf{Ref.} & \textbf{Departamento / Área} & \textbf{Efectivos} & \textbf{Porcentaje (\%)} \\
        \midrule
        \textbf{A} & Dirección y Gerencia General & 2 & 6.45\% \\
        \textbf{B} & Departamento de Producción & 10 & 32.26\% \\
        \textbf{C} & Ingeniería y Mantenimiento & 4 & 12.90\% \\
        \textbf{D} & Calidad e I+D+i (Laboratorio y Regulatorio) & 5 & 16.13\% \\
        \textbf{E} & Seguridad, Salud y Medio Ambiente (HSE) & 2 & 6.45\% \\
        \textbf{F} & Compras, Aprovisionamiento Agrícola y Logística & 3 & 9.68\% \\
        \textbf{G} & Administración, Finanzas, RRHH y Comercial & 5 & 16.13\% \\
        \midrule
        \multicolumn{2}{r}{\textbf{TOTAL TRABAJADORES:}} & \textbf{31} & \textbf{100.0\%} \\
        \bottomrule
    \end{tabular}
\end{table}

\section{Definición de Puestos, Formación y Funciones}

\subsection{A. Dirección y Gerencia (2 personas)}

\begin{itemize}[leftmargin=*]
    \item \textbf{1. Director/a General (CEO) (1 efectivo):}
          \begin{itemize}
              \item \textbf{Formación requerida:} Titulación universitaria superior en Ingeniería Química, Ingeniería Industrial, Biotecnología, Farmacia o afines, complementada con formación de posgrado en Administración y Dirección de Empresas (MBA o máster en gestión).
              \item \textbf{Funciones principales:}
                    \begin{itemize}
                        \item Definición de la estrategia corporativa y plan de negocio a corto, medio y largo plazo.
                        \item Gestión económica, financiera e interlocución con accionistas e inversores.
                        \item Relaciones institucionales de alto nivel con clientes clave, proveedores estratégicos y organismos públicos.
                        \item Supervisión general del rendimiento de todos los departamentos de la compañía.
                        \item Aprobación de inversiones de capital (CAPEX) y presupuestos operativos (OPEX).
                    \end{itemize}
          \end{itemize}

    \item \textbf{2. Director/a de Planta / Gerente Industrial (1 efectivo):}
          \begin{itemize}
              \item \textbf{Formación requerida:} Grado o Máster en Ingeniería Química o Ingeniería Industrial, con sólida experiencia contrastada ($\geq 5$ años) en gestión de plantas químicas, biotecnológicas o agroalimentarias.
              \item \textbf{Funciones principales:}
                    \begin{itemize}
                        \item Dirección integral de las operaciones fabriles: Producción, Mantenimiento, Calidad y Seguridad.
                        \item Coordinación técnica de la planta y optimización de costes de fabricación por kilogramo de producto.
                        \item Gestión de proyectos de ampliación, revamping y modificaciones de líneas de proceso.
                        \item Garantizar la continuidad operativa y la consecución del plan maestro de producción.
                    \end{itemize}
          \end{itemize}
\end{itemize}

\subsection{B. Departamento de Producción (10 personas)}

Este departamento constituye el núcleo operativo de la instalación, organizado para cubrir la operación continua del proceso de extracción y purificación.

\begin{itemize}[leftmargin=*]
    \item \textbf{1. Jefe/a de Producción (1 efectivo):}
          \begin{itemize}
              \item \textbf{Formación requerida:} Ingeniería Química o Ingeniería Industrial con experiencia en procesos industriales continuos/discontinuos.
              \item \textbf{Funciones:} Planificación de campañas de procesado, programación y gestión de turnos, seguimiento de rendimientos de extracción, control de consumos energéticos y materias auxiliares, y coordinación general del equipo operativo.
          \end{itemize}

    \item \textbf{2. Supervisores de Turno (2 efectivos):}
          \begin{itemize}
              \item \textbf{Formación requerida:} Grado en Ingeniería Química/Industrial o Técnico Superior en Procesos Químicos/Química Industrial (CFGS).
              \item \textbf{Funciones:} Supervisión directa de planta durante el turno de trabajo, control de variables de operación (temperaturas, caudales, presiones, pH, pureza), gestión de incidencias operativas inmediatas, garantía de cumplimiento estricto de POEs (Procedimientos Operativos Estandarizados) y coordinación de operadores.
          \end{itemize}

    \item \textbf{3. Operadores de Planta (6 efectivos):}
          \begin{itemize}
              \item \textbf{Formación requerida:} Ciclo Formativo de Grado Superior (CFGS) en Química Industrial, Operaciones de Planta Química, Fabricación de Productos Farmacéuticos o similar.
              \item \textbf{Funciones:}
                    \begin{itemize}
                        \item Operación de equipos y líneas de acondicionamiento de materia prima (molienda y extracción sólido-líquido).
                        \item Control de las etapas de clarificación, filtración tangencial, columnas de resinas de intercambio iónico y adsorción.
                        \item Operación de sistemas de concentración por evaporación al vacío, cristalización fraccionada y secado por atomización (\textit{spray drying}).
                        \item Tareas de envasado, muestreo rutinario y protocolos de limpieza y desinfección (\textit{CIP/SIP}).
                    \end{itemize}
          \end{itemize}

    \item \textbf{4. Responsable de Almacén y Silos (1 efectivo):}
          \begin{itemize}
              \item \textbf{Formación requerida:} CFGS en Gestión de Transporte y Logística o experiencia equivalente en almacenes industriales.
              \item \textbf{Funciones:} Recepción, pesaje y control de hojas secas de stevia en fardos/silos, gestión de reactivos químicos y disolventes, almacenamiento adecuado del producto terminado bajo condiciones controladas y preparación de expediciones.
          \end{itemize}
\end{itemize}

\subsection{C. Departamento de Ingeniería y Mantenimiento (4 personas)}

Garantiza la máxima fiabilidad de las instalaciones mecánicas, eléctricas y térmicas, así como la optimización continua del proceso.

\begin{itemize}[leftmargin=*]
    \item \textbf{1. Ingeniero/a de Procesos (1 efectivo):}
          \begin{itemize}
              \item \textbf{Formación requerida:} Grado/Máster en Ingeniería Química.
              \item \textbf{Funciones:} Elaboración y ajuste de balances de materia y energía, simulación de procesos (Aspen Plus / SuperPro Designer), dimensionamiento y optimización de reactores y columnas de resina, maximización del rendimiento de recuperación de rebaudiósido A, escalado y modificaciones de diagramas PFD/P\&ID.
          \end{itemize}

    \item \textbf{2. Ingeniero/a de Proyectos y Oficina Técnica (1 efectivo):}
          \begin{itemize}
              \item \textbf{Formación requerida:} Grado/Máster en Ingeniería Química o Ingeniería Industrial (Mecánica/Eléctrica).
              \item \textbf{Funciones:} Gestión de proyectos de nuevas instalaciones y modificaciones técnicas, elaboración y actualización de P\&ID y planos de implantación (\textit{layout}), especificación técnica de nuevos equipos (bombas, centrífugas, columnas, tanques), coordinación de contratistas y control de plazos y costes de proyecto.
          \end{itemize}

    \item \textbf{3. Técnicos de Mantenimiento Electromecánico (2 efectivos):}
          \begin{itemize}
              \item \textbf{Formación requerida:} CFGS en Mantenimiento Electromecánico, Mecatrónica Industrial o Mantenimiento de Instalaciones Térmicas y de Fluidos.
              \item \textbf{Funciones:} Ejecución del plan de mantenimiento preventivo y resolución de averías correctivas; mantenimiento de bombas centrífugas y de desplazamiento positivo, válvulas automáticas, intercambiadores de calor de placas/tubulares, motores eléctricos, lazos de instrumentación (sensores pH, caudalímetros, transmisores de presión) y servicios auxiliares (generador de vapor, aire comprimido, planta de agua osmotizada).
          \end{itemize}
\end{itemize}

\subsection{D. Departamento de Calidad e I+D+i (5 personas)}

Área crítica debido a que el producto es un edulcorante de alta pureza destinado a consumo humano bajo estrictas farmacopeas y regulaciones alimentarias (EFSA, FDA, JECFA).

\begin{itemize}[leftmargin=*]
    \item \textbf{1. Responsable de Calidad y Seguridad Alimentaria (1 efectivo):}
          \begin{itemize}
              \item \textbf{Formación requerida:} Licenciatura o Grado en Química, Bioquímica, Ciencia y Tecnología de los Alimentos o Ingeniería Química; formación acreditada en auditorías y normativas BRC, IFS, FSSC 22000 e ISO 9001.
              \item \textbf{Funciones:} Implantación y liderazgo del Sistema de Gestión de Calidad (SGC), plan APPCC/HACCP, homologación de proveedores de materia prima, gestión del protocolo de liberación de lotes, tratamiento de no conformidades, trazabilidad integral y atención de auditorías externas.
          \end{itemize}

    \item \textbf{2. Técnicos de Laboratorio de Control de Calidad (2 efectivos):}
          \begin{itemize}
              \item \textbf{Formación requerida:} CFGS en Laboratorio de Análisis y Control de Calidad, Grado en Química o Ingeniería Química.
              \item \textbf{Funciones:} Análisis de muestras de hojas recibidas (humedad, impurezas, contenido de esteviósido y Reb A), control analítico en proceso y control de producto final; cuantificación analítica mediante cromatografía líquida de alta resolución (\textbf{HPLC-UV/RID}), análisis de metales pesados, cenizas y recuentos microbiológicos.
          \end{itemize}

    \item \textbf{3. Ingeniero/a de I+D+i (1 efectivo):}
          \begin{itemize}
              \item \textbf{Formación requerida:} Grado/Máster en Ingeniería Química, Biotecnología o Química.
              \item \textbf{Funciones:} Investigación en la optimización de rendimientos de extracción con solventes verdes, formulación de nuevas mezclas de glucósidos enriquecidos en Rebaudiósido M y D, purificación enzimática, desarrollo de nuevos formatos (líquido, polvo granulado) y colaboración con centros tecnológicos y universidades.
          \end{itemize}

    \item \textbf{4. Técnico de Asuntos Regulatorios y Documentación (1 efectivo):}
          \begin{itemize}
              \item \textbf{Formación requerida:} Grado en Química, Tecnología de los Alimentos, Farmacia o Ingeniería Química.
              \item \textbf{Funciones:} Elaboración y actualización de especificaciones técnicas de producto, fichas de datos de seguridad (FDS), expedientes regulatorios para registro alimentario internacional, gestión de certificados Kosher, Halal y ecológico/vegano, y archivo documental técnico.
          \end{itemize}
\end{itemize}

\subsection{E. Seguridad, Salud Laboral y Medio Ambiente - HSE (2 personas)}

\begin{itemize}[leftmargin=*]
    \item \textbf{1. Responsable de HSE / PRL (1 efectivo):}
          \begin{itemize}
              \item \textbf{Formación requerida:} Grado en Ingeniería Química o Industrial, con Máster Oficial en Prevención de Riesgos Laborales (3 especialidades: Seguridad en el Trabajo, Higiene Industrial y Ergonomía/Psicosociología).
              \item \textbf{Funciones:} Evaluación de riesgos laborales, implantación del plan de prevención, formación en seguridad a trabajadores, planes de autoprotección y emergencias, seguridad de procesos químicos (ATEX por polvos y vapores de etanol) y control de sustancias peligrosas.
          \end{itemize}

    \item \textbf{2. Técnico de Medio Ambiente y Sostenibilidad (1 efectivo):}
          \begin{itemize}
              \item \textbf{Formación requerida:} Grado en Ciencias Ambientales, Ingeniería Química o Ingeniería Ambiental.
              \item \textbf{Funciones:} Gestión integral de residuos (especialmente orujo/torta de hoja agotada para compostaje/biomasa), gestión de la depuradora de aguas residuales industriales (EDARI), control de emisiones atmosféricas y vertidos, auditorías de huella hídrica y de carbono, y cumplimiento de la Autorización Ambiental Integrada (AAI).
          \end{itemize}
\end{itemize}

\subsection{F. Compras, Aprovisionamiento Agrícola y Logística (3 personas)}

\begin{itemize}[leftmargin=*]
    \item \textbf{1. Responsable de Compras Generales (1 efectivo):}
          \begin{itemize}
              \item \textbf{Funciones:} Negociación y aprovisionamiento de reactivos químicos, coadyuvantes de filtración, carbón activo, envases de grado alimentario, material de laboratorio y repuestos electromecánicos.
          \end{itemize}

    \item \textbf{2. Responsable de Aprovisionamiento Agrícola (1 efectivo - Puesto Estratégico):}
          \begin{itemize}
              \item \textbf{Formación requerida:} Grado en Ingeniería Agrícola / Agroalimentaria o experiencia en gestión de cultivos industriales.
              \item \textbf{Funciones:} Gestión de contratos agrarios directos con agricultores y cooperativas productoras de stevia, planificación de siembras y cosechas escalonadas, asesoramiento técnico agronómico (variedades ricas en Reb A), estimación de rendimientos por hectárea y control previo a cosecha.
          \end{itemize}

    \item \textbf{3. Técnico de Logística y Transporte (1 efectivo):}
          \begin{itemize}
              \item \textbf{Funciones:} Planificación y contratación de rutas de transporte nacional e internacional (contenedores marítimos/terrestres), gestión aduanera de exportación, coordinación de entregas a clientes industriales (bebidas, confitería, farmacia).
          \end{itemize}
\end{itemize}

\subsection{G. Administración, Finanzas, Recursos Humanos y Comercial (5 personas)}

\begin{itemize}[leftmargin=*]
    \item \textbf{1. Administración y Finanzas (2 efectivos):}
          \begin{itemize}
              \item \textbf{Formación requerida:} Grado en ADE, Economía, Finanzas o similar.
              \item \textbf{Funciones:} Contabilidad general y analítica de costes, facturación a clientes, gestión de cobros y pagos, control presupuestario mensual, tesorería y liquidación de impuestos.
          \end{itemize}

    \item \textbf{2. Técnico/a de Recursos Humanos (1 efectivo):}
          \begin{itemize}
              \item \textbf{Formación requerida:} Grado en Relaciones Laborales, Recursos Humanos, Psicología del Trabajo o ADE.
              \item \textbf{Funciones:} Gestión de contratación, nóminas, control horario, planes de acogida, evaluación del desempeño, plan anual de formación y coordinación con el servicio de prevención.
          \end{itemize}

    \item \textbf{3. Equipo Comercial / Ventas B2B (2 efectivos):}
          \begin{itemize}
              \item \textbf{Formación requerida:} Formación en Comercio, ADE o Ingeniería Agroalimentaria/Química con perfil técnico-comercial.
              \item \textbf{Funciones:} Prospección de mercado B2B (fabricantes de refrescos, lácteos, repostería, nutracéuticos), negociación de acuerdos de suministro a largo plazo, asistencia a ferias internacionales de ingredientes alimentarios y fidelización de clientes.
          \end{itemize}
\end{itemize}

\section{Organigrama de la Organización}

A continuación se presenta la estructura jerárquica y funcional de la planta industrial:

\begin{figure}[H]
\centering
\begin{tikzpicture}[
    level 1/.style={sibling distance=7.5cm, level distance=1.6cm},
    level 2/.style={sibling distance=2.25cm, level distance=2.3cm},
    box/.style={
        rectangle, 
        rounded corners=4pt, 
        draw=blue!80!black, 
        thick, 
        fill=blue!5, 
        align=center, 
        font=\small\bfseries, 
        inner sep=5pt,
        drop shadow={opacity=0.2}
    },
    dirbox/.style={
        rectangle, 
        rounded corners=5pt, 
        draw=red!80!black, 
        very thick, 
        fill=red!10, 
        align=center, 
        font=\normalsize\bfseries, 
        inner sep=6pt,
        drop shadow={opacity=0.3}
    },
    plantabox/.style={
        rectangle, 
        rounded corners=4pt, 
        draw=purple!80!black, 
        thick, 
        fill=purple!10, 
        align=center, 
        font=\small\bfseries, 
        inner sep=5pt,
        drop shadow={opacity=0.25}
    },
    deptbox/.style={
        rectangle, 
        rounded corners=3pt, 
        draw=teal!80!black, 
        thick, 
        fill=teal!10, 
        align=center, 
        font=\footnotesize\bfseries, 
        inner sep=4pt,
        text width=2.1cm,
        drop shadow={opacity=0.15}
    },
    edge from parent/.style={draw=black!70, thick, -{Latex[scale=0.8]}}
]

% Nodo Raíz: Dirección General
\node[dirbox] (DG) {Director/a General (CEO)\\\textmd{\footnotesize (1 persona)}}
    child { 
        node[box, fill=orange!10, draw=orange!80!black] (ADM) {Admón, RRHH y Comercial\\\textmd{\footnotesize (5 personas)}}
        child { node[deptbox] {Administración y Finanzas\\\textmd{\tiny (2)}} }
        child { node[deptbox] {Recursos Humanos\\\textmd{\tiny (1)}} }
        child { node[deptbox] {Ventas B2B y Comercial\\\textmd{\tiny (2)}} }
    }
    child { 
        node[plantabox] (DP) {Director/a de Planta\\\textmd{\footnotesize (1 persona)}}
        child { 
            node[deptbox] {Producción\\\textmd{\tiny (10 pers.)}} 
        }
        child { 
            node[deptbox] {Ingeniería y Mantenimiento\\\textmd{\tiny (4 pers.)}} 
        }
        child { 
            node[deptbox] {Calidad e I+D+i\\\textmd{\tiny (5 pers.)}} 
        }
        child { 
            node[deptbox] {Seguridad y Medio Ambiente\\\textmd{\tiny (2 pers.)}} 
        }
        child { 
            node[deptbox] {Compras y Aprov. Agrícola\\\textmd{\tiny (3 pers.)}} 
        }
    };

\end{tikzpicture}
\caption{Organigrama estructural por niveles de responsabilidad}
\label{fig:organigrama_general}
\end{figure}

\subsection{Detalle de Dependencias Internas del Área Industrial}

\begin{figure}[H]
\centering
\begin{tikzpicture}[
    node distance=1.1cm and 0.8cm,
    block/.style={rectangle, draw=blue!70!black, fill=blue!8, rounded corners=3pt, align=center, font=\scriptsize\bfseries, inner sep=4pt, text width=3.3cm},
    head/.style={rectangle, draw=purple!80!black, fill=purple!15, rounded corners=4pt, align=center, font=\small\bfseries, inner sep=5pt, text width=5cm},
    line/.style={draw=black!70, thick, -{Latex[scale=0.75]}}
]

\node[head] (dir) {Director/a de Planta / Gerente Industrial (1)};

% Nivel Jefaturas
\node[block, below left=1.2cm and 2.5cm of dir, fill=green!15, draw=green!60!black] (jprod) {Jefe de Producción (1)};
\node[block, below left=1.2cm and -0.8cm of dir, fill=blue!15, draw=blue!60!black] (jing) {Ing. Procesos (1) / Proyectos (1)};
\node[block, below right=1.2cm and -0.8cm of dir, fill=teal!15, draw=teal!60!black] (jcal) {Resp. Calidad (1) / I+D (1)};
\node[block, below right=1.2cm and 2.5cm of dir, fill=red!15, draw=red!60!black] (jhse) {Resp. HSE / PRL (1)};

% Nivel Subordinados
\node[block, below=0.6cm of jprod, font=\tiny\textmd] (subprod) {$\bullet$ 2 Supervisores Turno\\$\bullet$ 6 Operadores Planta\\$\bullet$ 1 Resp. Almacén};
\node[block, below=0.6cm of jing, font=\tiny\textmd] (subing) {$\bullet$ 2 Técnicos Mantenimiento Electromecánico};
\node[block, below=0.6cm of jcal, font=\tiny\textmd] (subcal) {$\bullet$ 2 Técnicos Lab. HPLC\\$\bullet$ 1 Técnico Regulatorio};
\node[block, below=0.6cm of jhse, font=\tiny\textmd] (subhse) {$\bullet$ 1 Técnico Medioambiental\\$\bullet$ Aprovisionamiento Agrícola (1)};

% Conexiones
\draw[line] (dir.south) -| (jprod.north);
\draw[line] (dir.south) -| (jing.north);
\draw[line] (dir.south) -| (jcal.north);
\draw[line] (dir.south) -| (jhse.north);

\draw[line] (jprod) -- (subprod);
\draw[line] (jing) -- (subing);
\draw[line] (jcal) -- (subcal);
\draw[line] (jhse) -- (subhse);

\end{tikzpicture}
\caption{Desglose funcional de dependencias técnicas en planta}
\label{fig:desglose_tecnico}
\end{figure}

\section{Conclusiones y Justificación Ingenieril del Modelo}

\begin{enumerate}
    \item \textbf{Adecuación a los requerimientos del proceso:} La tecnología de obtención de edulcorantes a partir de stevia requiere operaciones unitarias delicadas (cromatografía/resinas, filtración de membranas y cristalización). La asignación de 10 personas en producción y 4 en ingeniería/mantenimiento garantiza una operación continua y con un OEE (\textit{Overall Equipment Effectiveness}) elevado.
    \item \textbf{Garantía de calidad e inocuidad alimentaria:} Al ser un ingrediente alimentario de consumo directo o formulación, la dotación de 5 especialistas en Calidad, Laboratorio HPLC e I+D asegura el cumplimiento de las farmacopeas y normativas internacionales (pureza mínima del 95\% en glucósidos totales y ratio Reb A / Esteviósido controlado).
    \item \textbf{Eslabón agrícola estratégico:} La figura del \textit{Responsable de Aprovisionamiento Agrícola} es un factor diferenciador clave para mitigar la estacionalidad y variabilidad de la materia prima, asegurando contratos estables con agricultores locales y control de calidad en origen.
    \item \textbf{Sostenibilidad y seguridad integrada:} La inclusión de personal dedicado en exclusiva a HSE y Medio Ambiente responde a la necesidad de gestionar de forma circular los subproductos (orujo de hoja para compostaje o biomasa energética) y garantizar el control en zonas con riesgo ATEX por manipulación de polvos finos y disolventes.
\end{enumerate}